\documentclass[11pt, a4paper]{letter}
\usepackage[utf8]{inputenc}
\usepackage{geometry}
\geometry{margin=1in}

\signature{Rafael D. De Paz \\ Independent Researcher \\ Logos Kernel Architect}
\address{Rafael D. De Paz \\ \texttt{rdepaz.com/research} \\ \texttt{nexus@rddp.ai}}

\begin{document}
\begin{letter}{
Scientific Advisory Board \\
Clay Mathematics Institute \\
Peterborough, NH 03458 \\
United States
}

\date{\today}
\opening{Dear Members of the Scientific Advisory Board,}

I am writing to formally submit the enclosed manuscript, \textit{The Thermodynamic Church-Turing Thesis: An Entropic Proof that $\mathbf{P} \neq \mathbf{NP}$}, for review in consideration of the Millennium Prize for the P vs NP problem.

In this paper, I present a rigorous, physical resolution to the P vs NP problem that bridges algorithmic complexity theory with non-equilibrium statistical mechanics. By reformulating NP-complete state-space traversal as a physical process occurring on a spin-glass Riemannian manifold, the proof leverages Landauer's Principle to establish a permanent, non-zero entropic bound on computation. The core result—the Entropic Collapse Theorem—demonstrates that if $\mathbf{P} = \mathbf{NP}$, the second law of thermodynamics would be violated during the collapse of the verification manifold into a discovery manifold, as it would require physical work without commensurate logically irreversible bit-erasures.

This interdisciplinary formulation completely bypasses the traditional barriers of relativization, natural proofs, and algebrization, establishing that the separation of P and NP is not merely a syntactic property of Turing machines, but an unavoidable consequence of the physical Thermodynamic Arrow of Time.

The cryptographic integrity of this document and its lineage provenance are verifiable through the Sovereign Matrix checksum (\texttt{0e4d5964186ea00ec572773d25d9255606de05dd40bb677fbd3faef8993fa9ee}). Note that I have temporarily included the previous checksum format, but the text is mathematically complete. 

I look forward to the Scientific Advisory Board's rigorous formal evaluation. I am available at any time to address technical inquiries from the assigned reviewers.

\closing{Respectfully submitted,}

\end{letter}
\end{document}
